% \iffalse meta-comment
%
% Copyright (C) 2015-2016 by Zeping Lee <zepinglee@gmail.com>
%
% This file may be distributed and/or modified under the
% conditions of the LaTeX Project Public License, either
% version 1.2 of this license or (at your option) any later
% version. The latest version of this license is in:
%
%     http://www.latex-project.org/lppl.txt
%
% and version 1.2 or later is part of all distributions of
% LaTeX version 1999/12/01 or later.
%
% \fi
%
% \iffalse
%<*driver>
\ProvidesFile{ustcthesis.dtx}
%</driver>
%<class>\NeedsTeXFormat{LaTeX2e}[1999/12/01]
%<class>\ProvidesClass{ustcthesis}
%<*class>
    [2015/06/28 v2.0 Thesis Template for University of Science and Technology of China]
%</class>
%
%<*driver>
\documentclass[a4paper]{ltxdoc}
\EnableCrossrefs
\CodelineIndex
\RecordChanges

\usepackage{hyperref}
\hypersetup{
    CJKbookmarks=true,
    bookmarksnumbered=true,
    bookmarksopen=true,
    colorlinks=true,
    linkcolor=blue,
    citecolor=blue}
\usepackage[UTF8, heading, sub3section]{ctex}
\ctexset{
  abstractname   = 简介,
  indexname      = 代码索引,
  section        = {
    format = \Large\bfseries\raggedright,
    name   = {第,节},
  },
}
\usepackage{tabularx}
\usepackage{booktabs}


\begin{document}
    \DocInput{ustcthesis.dtx}
\end{document}
\usepackage{metalogo}
\usepackage{doc}
%</driver>
% \fi
%
% \CheckSum{0}
%
% \CharacterTable
%  {Upper-case    \A\B\C\D\E\F\G\H\I\J\K\L\M\N\O\P\Q\R\S\T\U\V\W\X\Y\Z
%   Lower-case    \a\b\c\d\e\f\g\h\i\j\k\l\m\n\o\p\q\r\s\t\u\v\w\x\y\z
%   Digits        \0\1\2\3\4\5\6\7\8\9
%   Exclamation   \!     Double quote  \"     Hash (number) \#
%   Dollar        \$     Percent       \%     Ampersand     \&
%   Acute accent  \'     Left paren    \(     Right paren   \)
%   Asterisk      \*     Plus          \+     Comma         \,
%   Minus         \-     Point         \.     Solidus       \/
%   Colon         \:     Semicolon     \;     Less than     \<
%   Equals        \=     Greater than  \>     Question mark \?
%   Commercial at \@     Left bracket  \[     Backslash     \\
%   Right bracket \]     Circumflex    \^     Underscore    \_
%   Grave accent  \`     Left brace    \{     Vertical bar  \|
%   Right brace   \}     Tilde         \~}
%
% \changes{v2.0}{2015/06/26}{Initial Build}
%
% \GetFileInfo{ustcthesis.dtx}
%
% \DoNotIndex{\newcommand,\newenvironment}
%
% \title{\textsf{ustcthesis}使用说明}
% \author{seisman\thanks{seisman.info@gmail.com} \and
%         Zeping Lee\thanks{zepinglee@gmail.com}}
%
% \maketitle
%
% \section{简介}
%
% \textsf{ustcthesis}是用来排版中国科学技术大学本科和研究生毕业论文的\LaTeX{}模板，
% 按照《中国科学技术大学研究生学位论文撰写规范》
% \footnote{http://gradschool.ustc.edu.cn/ylb/material/xw/wdxz/1.doc}
% 和《关于本科毕业论文（设计）格式和统一封面的通知》
% \footnote{http://www.teach.ustc.edu.cn/document/doc-administration/4032.html}
% 的要求编写。
%
% 其前身是中国科学技术大学本科论文模板（作者XPS，最后维护ywg）和中国科学技术大学研究生论文
% 模板（作者Liuqs，主要维护Liuqs、Guolicai）。
% 后来两模板进行了整合梳理，由ywg维护。
% 2015年，seisman和zepinglee基于\textsf{ctex} 2.0重新编写了模板。\pkg{ctex}
%
% \section{使用说明}
%
% \subsection{文件组成}
% 本模板的主要文件如表\ref{tab:file}：
%
% \begin{table}[htbp]
% \centering
% \begin{tabularx}{\linewidth}{llX}
%   \toprule
%   类别     & 文件                    & 说明 \\
%   \midrule
%   模板文件 & \texttt{ustcthesis.dtx} & 主要目的模板文件，可生成
%                            \texttt{ustcthesis.cls}和\texttt{ustcthesis.pdf} \\
%           & \texttt{ustcthesis.bst} & 参考文献列表的格式文件 \\
%   \midrule
%   生成文件 & \texttt{ustcthesis.pdf} & （你正在阅读的）模板使用说明 \\
%           & \texttt{ustcthesis.pdf} & 文档类文件 \\
%   \midrule
%   示例文件 & \texttt{main.tex}       & 主文档 \\
%           & \texttt{chapters/*.tex} & 示例文档的各个章节 \\
%           & \texttt{figures/}       & 放置图片的目录 \\
%           & \texttt{bib/tex.bib}    & 参考文献示例数据库 \\
%   \midrule
%   其他    & \texttt{Makefile}       & 用于GNU make \\
%   \bottomrule
% \end{tabularx}
% \caption{\textsf{ustcthesis}的文件组成}
% \label{tab:file}
% \end{table}
%
% \section{排版注意事项}
% \begin{itemize}
%   \item 英文逗号、句号、问号和叹号后面要加空格。
% \end{itemize}
%
% \StopEventually{\PrintIndex}
%
% \linespread{1.2}\selectfont
% \section{代码实现}
% \subsection{声明选项}
%    \begin{macrocode}
\RequirePackage{kvoptions}
\SetupKeyvalOptions{family=ustc@opt, prefix=ustc@opt@, setkeys=\kvsetkeys}
\DeclareBoolOption[false]{doctor}
\DeclareBoolOption[false]{master}
\DeclareBoolOption[false]{bachelor}
\DeclareBoolOption[false]{pdf}
\DeclareComplementaryOption{print}{pdf}
%    \end{macrocode}
% |ustc@page|这一family用来设置分页方式
%    \begin{macrocode}
\SetupKeyvalOptions{family=ustc@page, prefix=ustc@page@, setkeys=\kvsetkeys}
\DeclareBoolOption[false]{oneside}
\DeclareComplementaryOption{twoside}{oneside}
\DeclareBoolOption[true]{openany}
\DeclareComplementaryOption{openright}{openany}
\DeclareBoolOption[true]{pageright}
\DeclareComplementaryOption{pagecenter}{pageright}
\DeclareDefaultOption{\PassOptionsToClass{\CurrentOption}{ctexbook}}
%    \end{macrocode}
% 先处理ustc@opt这一family，然后设置本科和pdf的默认分页方式为oneside, pageright
%    \begin{macrocode}
\ProcessKeyvalOptions{ustc@opt}
\ifustc@opt@bachelor
    \ustc@page@onesidetrue
    \ustc@page@pagerightfalse
\fi
\ifustc@opt@pdf
    \ustc@page@onesidetrue
    \ustc@page@pagerightfalse
\fi
%    \end{macrocode}
% 处理其他的选项
%    \begin{macrocode}
\ProcessKeyvalOptions*\relax
\ifustc@page@oneside
    \PassOptionsToClass{oneside}{ctexbook}
\fi
\ifustc@page@openany
    \PassOptionsToClass{openany}{ctexbook}
\fi
%    \end{macrocode}
% 默认情况下\textsf{fontspec}会将|\mathrm{}|的字体更改为Times New Roman，
% 而|\mathnormal|依然为Computer Modern，比Times New Roman要细，
% 所以设置\textsf{fontspec}不处理数学字体
%    \begin{macrocode}
\PassOptionsToPackage{no-math}{fontspec}
%    \end{macrocode}
% 载入\textsf{ctexbook}文档类，注意要求为ctex 2.0以及更高的版本。
%    \begin{macrocode}
\LoadClass[a4paper, UTF8, zihao=-4]{ctexbook}[2015/06/19]
%    \end{macrocode}
%
% \subsection{设置西文字体}
% 设置西文字体，中文字体已经由ctex自动设置，
%    \begin{macrocode}
\setmainfont{Times New Roman}
\setsansfont{Arial}
\setmonofont{Courier New}
%    \end{macrocode}
%
% \subsection{超链接与书签}
%    \begin{macrocode}
\RequirePackage{hyperref}
\hypersetup{
    CJKbookmarks=true,
    bookmarksnumbered=true,
    bookmarksopen=true,
    pdfborder=0 0 0}
%    \end{macrocode}
% 如果为pdf样式，设置hyperlink颜色
%    \begin{macrocode}
\ifustc@opt@pdf
    \hypersetup{
        colorlinks=true,
        linkcolor=blue,
        citecolor=blue}
\fi
%    \end{macrocode}
% 设置pdf的信息
%    \begin{macrocode}
\AtBeginDocument{
    \hypersetup{
        pdftitle={\ustc@title},
        pdfauthor={\ustc@author}
    }
}
%    \end{macrocode}
%
% \subsection{定义命令}
% \begin{macro}{\setfontsize}
% 直接设置字体大小的命令，可选设置行距，默认单倍行距
%    \begin{macrocode}
\RequirePackage{xparse}
\DeclareDocumentCommand \setfontsize { m o } {
    \IfValueTF {#2}
        { \fontsize{#1}{#2} }
        { \fontsize{#1}{1.2 \dimexpr#1} }
    \linespread{1}\selectfont\relax}
%    \end{macrocode}
% \end{macro}
% \begin{macro}{\zhspace}
% 生成整数倍汉字宽度的空格
%    \begin{macrocode}
\newcommand\zhspace[1][1]{\hspace{#1\ccwd}}
%    \end{macrocode}
% \end{macro}
% \begin{macro}{\ustc@underline}
% 生成空的下划线
%    \begin{macrocode}
\newcommand\ustc@underline[2][6em]{
    \hskip1pt\underline{\hb@xt@ #1{\hss#2\hss}}\hskip3pt
}
%    \end{macrocode}
% \end{macro}
% \begin{macro}{\ustc@define@term}
% 声明“作者姓名”等项目的命令
%    \begin{macrocode}
\def\ustc@define@term#1{
    \expandafter\gdef\csname #1\endcsname##1{
        \expandafter\gdef\csname ustc@#1\endcsname{##1}
    }
    \csname #1\endcsname{}
}
\ustc@define@term{title}
\ustc@define@term{secrettext}
\ustc@define@term{author}
\ustc@define@term{depart}
\ustc@define@term{major}
\ustc@define@term{advisor}
\ustc@define@term{coadvisor}
\ustc@define@term{studentid}
\ustc@define@term{submitdate}
\ustc@define@term{entitle}
\ustc@define@term{ensecrettext}
\ustc@define@term{enauthor}
\ustc@define@term{enmajor}
\ustc@define@term{enadvisor}
\ustc@define@term{encoadvisor}
\ustc@define@term{ensubmitdate}
%    \end{macrocode}
% \end{macro}
%
% \subsection{汉化}
%    \begin{macrocode}
\ctexset{
    contentsname = {目\hspace{\ccwd}录},
}
%    \end{macrocode}
%
% \subsection{页面设置}
% 首先用\textsf{geometry}宏包设置纸张和页面。
%    \begin{macrocode}
\RequirePackage{geometry}
\geometry{
    paper=a4paper,
    top=2.54cm, bottom=2.54cm,
    left=3.17cm, right=3.17cm,
    headsep=0.74cm, headheight=0cm,
    footskip=0.79cm}
%    \end{macrocode}
% 再用\textsf{titleps}宏包设置页眉页码，定义了front, main两个页面风格分别用于前言和正文。
% |\ifustc@page@pageright|控制是否将页码置于外侧。
%    \begin{macrocode}
\RequirePackage{titleps}
\newpagestyle{front}[\zihao{5}]{
    \sethead{}{\chaptertitle}{}
    \ifustc@page@pageright
        \setfoot[\thepage][][]{}{}{\thepage}
    \else
        \setfoot{}{\thepage}{}
    \fi
    \headrule}
\newpagestyle{main}[\zihao{5}]{
    \sethead{}{\CTEXthechapter\quad\chaptertitle}{}
    \ifustc@page@pageright
        \setfoot[\thepage][][]{}{}{\thepage}
    \else
        \setfoot{}{\thepage}{}
    \fi
    \headrule}
%    \end{macrocode}
% 重新定义|\frontmatter|, |\mainmatter|, |\backmatter|，
% 使得在恰当位置应用正确的页面风格。
%    \begin{macrocode}
\renewcommand\frontmatter{
    \cleardoublepage
    \@mainmatterfalse
    \pagenumbering{Roman}
    \pagestyle{front}
    \ctexset{chapter/pagestyle = front}}
\renewcommand\mainmatter{%
    \cleardoublepage
    \@mainmattertrue
    \pagenumbering{arabic}
    \pagestyle{main}
    \ctexset{chapter/pagestyle = main}}
\renewcommand\backmatter{%
    \if@openright
        \cleardoublepage
    \else
        \clearpage
    \fi
    \@mainmatterfalse
    \pagestyle{front}
    \ctexset{chapter/pagestyle = front}}
%    \end{macrocode}
%
% \subsection{扉页}
% \begin{macro}{\maketitle}
% 重新定义|\maketitle|，分别调用|\make@cntitle|, |\make@entitle|生成中、英文扉页。
%    \begin{macrocode}
\newcommand{\ustc@declaretext} {
本人声明所呈交的学位论文，是本人在导师指导下进行研究工作所取得的成果。
除已特别加以标注和致谢的地方外，论文中不包含任何他人已经发表或撰写过
的研究成果。与我一同工作的同志对本研究所做的贡献均已在论文中作了明确的说明。
}

\newcommand{\ustc@authorization}{
作为申请学位的条件之一，学位论文著作权拥有者授权中国科学技术大学拥有
学位论文的部分使用权，即：学校有权按有关规定向国家有关部门或机构送交
论文的复印件和电子版，允许论文被查阅和借阅，可以将学位论文编入《中国
学位论文全文数据库》等有关数据库进行检索，可以采用影印、缩印或扫描等
复制手段保存、汇编学位论文。本人提交的电子文档的内容和纸质论文的内容
相一致。

保密的学位论文在解密后也遵守此规定。
}

\newcommand\make@statement{
    \setfontsize{12bp}[20bp]
    \vspace*{0.5cm}
    \begin{center}
    \zihao{3}\textbf{中国科学技术大学学位论文原创性声明}
    \end{center}

    \par\ustc@declaretext\par

    \vskip 1cm
    作者签名： \ustc@underline[6em]{} \hspace{5em} 签字日期：\ustc@underline[8em]{}
    \vskip 3cm

    \begin{center}
    \zihao{3}\textbf{中国科学技术大学学位论文授权使用声明}
    \end{center}

    \par\ustc@authorization\par
    \vskip 0.5cm
    \raisebox{1.1mm}{\fbox{\rule{0pt}{1.4mm}\rule{1.4mm}{0pt}}}\ 公开
    \hspace{0.5cm}
    \raisebox{1.1mm}{\fbox{\rule{0pt}{1.4mm}\rule{1.4mm}{0pt}}}
    保密（\ustc@underline[2em]{}\hspace{0.1em} 年）

    \vskip 0.5cm

    {\renewcommand{\arraystretch}{2.0}
    \begin{tabular}{p{6cm}p{6cm}}
    作者签名：\ustc@underline{} & 导师签名：\ustc@underline{} \\
    签字日期：\ustc@underline{} & 签字日期：\ustc@underline{} \\
    \end{tabular}}
}

\renewcommand\maketitle{
    \newgeometry{
        top=3.8cm, bottom=3.8cm,
        left=3.2cm, right=3.2cm,
        headheight=0cm, headsep=0.8cm,
        footskip=0.8cm}
    \pagestyle{empty}
    \pdfbookmark[-1]{\ustc@title}{title}
    \make@cntitle \cleardoublepage
    \make@entitle \cleardoublepage
    \restoregeometry
    \ifustc@opt@bachelor\relax\else
        \make@statement \cleardoublepage
    \fi}
%    \end{macrocode}
% \end{macro}
%
% \begin{macro}{\make@cntitle}
% 扉页为了排版``距离顶边xx厘米''，使用了tabular嵌套parbox的方法。
% 如果使用|\vpace{}|会产生额外空隙；如果叠parbox也不行。
% 注意这里没有严格按照《规范》的格式范例的尺寸，其中``学位论文''向上调整了0.7cm，
% 论文标题和作者姓名部分分别在parbox内部进行了垂直居中处理。另外两个eps图片为手工测量所得。
%    \begin{macrocode}
\newcommand\make@cntitle{
    \begin{titlepage}
    \begin{center}
    \renewcommand{\arraystretch}{0}
    \begin{tabular}{@{}c@{}}
        \\[0.2cm]
        \parbox[t][1.2cm][t]{\textwidth}
            {\raggedleft\fangsong\setfontsize{14bp}\ustc@secrettext} \\
        \parbox[t][2.7cm][t]{\textwidth}{\centering
            \includegraphics[width=11cm]{figures/ustc_logo_text.eps}} \\
        \parbox[t][3cm][t]{\textwidth}{
            \centering\bfseries\setfontsize{56bp}
            \ifustc@opt@doctor 博士学位论文 \fi
            \ifustc@opt@master 硕士学位论文 \fi
            \ifustc@opt@bachelor 学士学位论文 \fi} \\
        \includegraphics[height=4.9cm]{figures/ustc_logo_fig.eps} \\
        \parbox[t][4.7cm][c]{\textwidth}
            {\centering\bfseries\setfontsize{26bp}\ustc@title} \\
        \parbox[t][4.2cm][c]{\textwidth}{\setfontsize{16bp}
            \hspace*{2.8cm}{\bfseries 作者姓名：}\hfill \ustc@author \hfill\mbox{} \\
            \hspace*{2.8cm}{\bfseries 学科专业：}\hfill \ustc@major \hfill\mbox{} \\
            \hspace*{2.8cm}{\bfseries 导师姓名：}\hfill \ustc@advisor \hfill\mbox{} \\
            \hspace*{2.8cm}{\bfseries 完成时间：}\hfill \ustc@submitdate \hfill\mbox{}}
    \end{tabular}
    \end{center}
    \end{titlepage}}
%    \end{macrocode}
% \end{macro}
% \begin{macro}{\make@entitle}
% 由于《规范》对宽度的限制，英文扉页 Supervisor 一栏放不下两个英文名称
%    \begin{macrocode}
\newcommand\make@entitle{
    \begin{titlepage}
    \begin{center}
    \renewcommand{\arraystretch}{0}
    \begin{tabular}{@{}c@{}}
        \\[0.2cm]
        \parbox[t][1cm][t]{\textwidth}
            {\raggedleft\setfontsize{14bp}\ustc@ensecrettext} \\
        \parbox[t][1cm][t]{\textwidth}{\centering\sffamily\setfontsize{20bp}[30bp]
            University of Science and Technology of China} \\
        \parbox[t][2.2cm][t]{\textwidth}{\centering\bfseries\setfontsize{26bp}[30bp]
            \ifustc@opt@doctor A dissertation for doctor's degree \fi
            \ifustc@opt@master A dissertation for master's degree \fi
            \ifustc@opt@bachelor A dissertation for bachelor's degree \fi} \\
        \includegraphics[height=4.9cm]{figures/ustc_logo_fig.eps} \\
        \\[2.9cm]
        \parbox[t][4.7cm][t]{\textwidth}
            {\centering\bfseries\setfontsize{26bp}[30bp]\ustc@entitle} \\
        \parbox[t][4cm][c]{\textwidth}{\setfontsize{16bp}[30bp]
            \hspace*{2.8cm} Author: \hspace{1.6cm}\hfill \ustc@enauthor \hfill\mbox{} \\
            \hspace*{2.8cm} Speciality: \hspace{0.8cm}\hfill \ustc@enmajor \hfill\mbox{} \\
            \hspace*{2.8cm} Supervisor: \hspace{0.8cm}\hfill \ustc@enadvisor \hfill\mbox{} \\
            \hspace*{2.8cm} Finished Time: \hfill \ustc@ensubmitdate \hfill\mbox{}}
    \end{tabular}
    \end{center}
    \end{titlepage}}
%    \end{macrocode}
% \end{macro}
%
% \subsection{摘要}
% \begin{environment}{cnabstract}
% 中文摘要环境
%    \begin{macrocode}
\newenvironment{cnabstract}{\chapter[摘要]{摘\hspace{\ccwd}要}}{\cleardoublepage}
%    \end{macrocode}
% \end{environment}
% \begin{environment}{enabstract}
% 英文摘要环境
%    \begin{macrocode}
\newenvironment{enabstract}{
    \ctexset { chapter/format += \sffamily\bfseries\setfontsize{16bp} }
    \chapter[Abstract]{ABSTRACT}}
    {\cleardoublepage}
%    \end{macrocode}
% \end{environment}
% 两个生成关键字的命令，其中|\phantom{}|是为了空一行，
%    \begin{macrocode}
\newcommand\keywords[1]{\par\phantom{关键词}\par\noindent\hangindent=4\ccwd\relax
    \textbf{关键词：}#1}
\newcommand\enkeywords[1]{\par\phantom{keywords}\par\noindent\hangindent=5.3em\relax
    \textbf{Key Words:} #1}
%    \end{macrocode}
%
% \subsection{目录}
% \begin{macro}{\tableofcontents}
% 调用\textsf{tocbibind}宏包将目录、参考文献加入目录中。
% 为了让目录中的“目录”二字中间没有空格，用\textsf{tocbibind}的方法重新定义|\tableofcontents|
%    \begin{macrocode}
\RequirePackage{tocbibind}
\renewcommand{\tableofcontents}{\tocfile{目录}{toc}}
%    \end{macrocode}
% \end{macro}
% 下面用\textsf{titletoc}宏包设置每级目录的格式
%    \begin{macrocode}
\RequirePackage{titletoc}
\newcommand\ustc@leaders{\titlerule*[0.5pc]{$\cdot$}} % 指引线
\titlecontents{chapter}
    [0bp]
    {\vspace{6bp}\zihao{4}}
    {\thecontentslabel\hspace*{0.5em}}
    {}{\ustc@leaders\contentspage}
\titlecontents{section}
    [\ccwd]
    {\vspace{6bp}\zihao{-4}}
    {\thecontentslabel\hspace*{0.5em}}
    {}{\ustc@leaders\contentspage}
\titlecontents{subsection}
    [2\ccwd]
    {\vspace{6bp}\zihao{5}}
    {\thecontentslabel\hspace*{0.5em}}
    {}{\ustc@leaders\contentspage}
%    \end{macrocode}
% 同时设置图、表的目录格式
%    \begin{macrocode}
\titlecontents{figure}
    [\ccwd]
    {\zihao{-4}}
    {\thecontentslabel\hspace*{0.5em}}
    {}{\ustc@leaders\contentspage}
\titlecontents{table}
    [\ccwd]
    {\zihao{-4}}
    {\thecontentslabel\hspace*{0.5em}}
    {}{\ustc@leaders\contentspage}
%    \end{macrocode}
%
% \subsection{正文}
% 用\textsf{ctex}的接口设置全局的章节标题格式，注意由于\textsf{ctex} 2.1以下版本的bug，
% chapter的nameformat和titleformat的这两行必须置空。
% 另外摘要、参考文献等章节也是用|\chapter{}|在底层实现的
%    \begin{macrocode}
\ctexset{
    chapter = {
        format += \bfseries\setfontsize{16bp},
        nameformat = {},
        titleformat = {},
        number = \arabic{chapter},
        beforeskip = 24bp,
        afterskip = 18bp,
    },
    section = {
        format += \raggedright\bfseries\setfontsize{14bp},
        beforeskip = 24bp,
        afterskip = 6bp,
    },
    subsection = {
        format+= \raggedright\bfseries\setfontsize{13bp},
        beforeskip = 12bp,
        afterskip = 6bp,
    },
    subsubsection = {
        format+= \raggedright\bfseries\setfontsize{12bp},
        beforeskip = 12bp,
        afterskip = 6bp,
    }
}
%    \end{macrocode}
% 正文章节编号到3级section标题
%    \begin{macrocode}
\setcounter{secnumdepth}{3}
%    \end{macrocode}
% 正文段落文字格式设置
%    \begin{macrocode}
\setfontsize{12bp}[20bp]
\renewcommand\normalsize{\setfontsize{12bp}[20bp]}
\setlength{\parskip}{0bp}
%    \end{macrocode}
%
% \subsection{图、表}
% 用\textsf{caption}宏包设置图、表的格式
%    \begin{macrocode}
\RequirePackage{caption}
\DeclareCaptionLabelSeparator{zhspace}{\hspace{\ccwd}}
\captionsetup{
    format = hang,
    labelfont = {bf, small},
    labelsep = zhspace,
    textfont = small}
\captionsetup[figure]{
    position = bottom,
    aboveskip = 6bp,
    belowskip = 12bp}
\captionsetup[table]{
    position = top,
    aboveskip = 6bp,
    belowskip = 6bp}
%    \end{macrocode}
% \begin{macro}{\note}
% 新定义了|\note{}|来生成图标的附注。
% 如果用|\caption{}|生成图标的附注会导致图标的序号有误；
% 如果用|\bicaption{}|会导致表注无法置于表后，而且对齐方式不对。
%    \begin{macrocode}
\newcommand\note[1]{
    \captionsetup{position = bottom}
    \caption*{\hangindent=2\ccwd\relax\textbf{注}：#1}}
%    \end{macrocode}
% \end{macro}
%
% \subsection{数学公式}
% \LaTeX{}默认10pt时有|\DeclareMathSizes{\@xpt}{\@xpt}{7}{5}|，
% 这里按照相同的比例设置 10.5bp。
%    \begin{macrocode}
\DeclareMathSizes{10.5bp}{10.5bp}{7.35bp}{5.25bp}
%    \end{macrocode}
%
% \subsection{符号说明}
% \begin{environment}{notation}
% 这里只定义了一个满足要求的环境，具体内容由用户自定义
%    \begin{macrocode}
\newenvironment{notation}{
    \ctexset{chapter/format += \rmfamily\setfontsize{12bp}}
    \chapter{符号说明}
    \setfontsize{10.5bp}[16bp]
    \setlength{\itemsep}{0bp}}{}
%    \end{macrocode}
% \end{environment}
%
% \subsection{参考文献}
% 目前\BibTeX 使用的.bst格式是南京大学胡海星编写的
% \footnote{https://github.com/Haixing-Hu/GBT7714-2005-BibTeX-Style}，
% 不支持author-year式的文献标注和列表，故采用GB/T 7714-2005规定的numeric风格。
%    \begin{macrocode}
\usepackage[super,sort&compress]{natbib}
\setcitestyle{square}
\bibliographystyle{ustcthesis}
\renewcommand\bibfont{\setfontsize{10.5bp}[20bp]}
\setlength{\bibsep}{0bp}
\setlength{\bibhang}{1em}
%    \end{macrocode}
% \begin{macro}{\bibliography}
% 由于|\bibliography{}|位于正文中，页眉会默认显示“附录 参考文献”，所以修正页面风格：
%    \begin{macrocode}
\let\ustc@save@bibliography=\bibliography
\renewcommand\bibliography[1]{%
    \if@openright
        \cleardoublepage
    \else
        \clearpage
    \fi
    \pagestyle{front}
    \ctexset{chapter/pagestyle = front}
    \ustc@save@bibliography{#1}}
%    \end{macrocode}
% \end{macro}
%
% \subsection{附录}
% \begin{macro}{\appendix}
% 由于|\bibilography{}|处更改了页面风格，必须将附录的再更改为main。
%    \begin{macrocode}
\let\ustc@save@appendix=\appendix
\renewcommand\appendix{%
    \if@openright
        \cleardoublepage
    \else
        \clearpage
    \fi
    \pagestyle{main}
    \ctexset{chapter/pagestyle = main}
    \ustc@save@appendix}
%    \end{macrocode}
% \end{macro}
% \begin{environment}{acknowledgements}
% 定义了一个满足要求的致谢环境：
%    \begin{macrocode}
\newenvironment{acknowledgements}{\chapter[致谢]{致\hspace{\ccwd}谢}}{}
%    \end{macrocode}
% \end{environment}
% \begin{environment}{acknowledgements}
% 发表成果环境：
%    \begin{macrocode}
\newenvironment{publications}{\chapter{在读期间发表的学术论文与取得的研究成果}}{}
%    \end{macrocode}
% \end{environment}
%    \begin{macrocode}
% 书脊
% 书脊的论文题目、系别、作者姓名默认为\ustc@title, \ustc@depart, \ustc@author，
% 竖排的内部实现是，先用 \addCJKfontfeatures{Vertical=RotatedGlyphs} 旋转每个字，
% 再用 \rotatebox{-90}{} 旋转整个盒子，
% 由于中文字形的旋转后的深度会变，所以使用 \raisebox 进行调整。
% 所以留下 \spinetitle 等接口进行设置，如果标题里含有英文，要用\entext{}包括起来
% \RequirePackage{rotating}
% \ustc@define@term{spinetitle}
% \ustc@define@term{spinedepart}
% \ustc@define@term{spineauthor}
%
% \spinetitle{\ustc@title}
% \spinedepart{\ustc@depart}
% \spineauthor{\ustc@author}
%
% \newcommand\entext[1]{\raisebox{-0.75ex}{~#1{}~}}
%
% \newcommand\make@spine{
%     \begin{titlepage}
%     \begin{center}
%         书脊(此页仅用于制作书脊,不用单独打印放入论文)\par
%         \renewcommand\\{\relax}
%         \sffamily\bfseries\setfontsize{12bp}[14bp]\ziju{0.166667}
%         \addCJKfontfeatures{Vertical=RotatedGlyphs}
%         \rotatebox{-90}{\fbox{\parbox{20cm}{
%             \hspace*{3cm plus 1cm minus 2cm}
%             \raisebox{0.75ex}{\ustc@spinetitle} \hfill\quad\hfill
%             \raisebox{0.75ex}{\ustc@spinedepart} \hfill\quad\hfill
%             \raisebox{0.75ex}{\ustc@spineauthor} \hfill\quad\hfill
%             \raisebox{0.75ex}{中国科学技术大学}
%             \hspace{3cm plus 1cm minus 2cm}
%             \vskip 5.8pt\relax}}}
%     \end{center}
%     \end{titlepage}
% }
%
% \spinetitle{\entext{USTC}本硕博毕业论文示例文档} % 书脊的英文文字要用 \entext{} 括起来

\endinput

%    \end{macrocode}
%
% \Finale
\endinput
